\documentclass[11pt, a4paper]{article}
\usepackage[utf8]{inputenc}
\usepackage[T1]{fontenc}
\usepackage{babel}
\babelprovide[import]{spanish}
\usepackage{geometry}
\usepackage{graphicx}
\usepackage{float}
\usepackage{amsmath, amssymb}
\usepackage{mathtools}
\usepackage{booktabs}
\usepackage{array}
\usepackage{xcolor}
\usepackage{colortbl}
\usepackage{enumitem}
\usepackage{listings}
\usepackage{caption}
\usepackage{hyperref}
\usepackage{fancyhdr}
\usepackage{titlesec}
\usepackage{tcolorbox}
\tcbuselibrary{skins, breakable}

\geometry{a4paper, left=2.4cm, right=2.4cm, top=2.8cm, bottom=2.8cm, headheight=15pt}

\definecolor{navy}{HTML}{1a3a5c}
\definecolor{rojo}{HTML}{c8392b}
\definecolor{verde}{HTML}{27ae60}
\definecolor{grisclaro}{HTML}{f4f6f8}
\definecolor{azulclaro}{HTML}{2980b9}
\definecolor{codebg}{HTML}{1e2736}
\definecolor{codefg}{HTML}{e8eaf0}
\definecolor{codekw}{HTML}{56b6c2}
\definecolor{codestr}{HTML}{98c379}
\definecolor{codecom}{HTML}{7f848e}
\definecolor{linenum}{HTML}{4b5263}

\hypersetup{colorlinks=true, linkcolor=navy, urlcolor=azulclaro, citecolor=rojo,
  pdftitle={SARIMA -- Tres Ejemplos con Datos Financieros},
  pdfauthor={Dr. Cesar Galindo -- Universidad Panamericana}}

\pagestyle{fancy}
\fancyhf{}
\fancyhead[L]{\small\color{navy}\textbf{Series de Tiempo para Finanzas}}
\fancyhead[R]{\small\color{navy}Universidad Panamericana $\cdot$ Verano 2026}
\fancyfoot[C]{\small\thepage}
\renewcommand{\headrulewidth}{0.5pt}
\renewcommand{\headrule}{\hbox to\headwidth{\color{navy}\leaders\hrule height \headrulewidth\hfill}}

\titleformat{\section}{\large\bfseries\color{navy}}{\thesection.}{0.6em}{}
  [\vspace{-4pt}\textcolor{navy}{\titlerule[1pt]}\vspace{4pt}]
\titleformat{\subsection}{\normalsize\bfseries\color{rojo}}{\thesubsection.}{0.5em}{}
\titleformat{\subsubsection}{\normalsize\bfseries\color{azulclaro}}{\thesubsubsection.}{0.5em}{}

\lstdefinestyle{python}{
  language=Python,
  backgroundcolor=\color{codebg},
  basicstyle=\ttfamily\scriptsize\color{codefg},
  keywordstyle=\bfseries\color{codekw},
  stringstyle=\color{codestr},
  commentstyle=\itshape\color{codecom},
  numberstyle=\tiny\color{linenum},
  numbers=left, numbersep=8pt, stepnumber=5,
  breaklines=true, breakatwhitespace=true,
  showstringspaces=false, tabsize=4,
  frame=single, framerule=0pt, rulecolor=\color{codebg},
  xleftmargin=2em, framexleftmargin=1.8em, captionpos=b,
}

\tcbset{
  teorema/.style={enhanced, breakable, colback=grisclaro, colframe=navy,
    fonttitle=\bfseries\color{white}, coltitle=white,
    attach boxed title to top left={yshift=-2mm, xshift=4mm},
    boxed title style={colback=navy}, top=6pt, bottom=6pt, left=6pt, right=6pt},
  formula/.style={enhanced, colback=navy!8, colframe=navy!60,
    fonttitle=\bfseries\small, top=6pt, bottom=6pt, left=8pt, right=8pt, arc=3pt},
  aplicacion/.style={enhanced, colback=verde!6, colframe=verde!50,
    fonttitle=\bfseries\small\color{verde!60!black},
    top=5pt, bottom=5pt, left=6pt, right=6pt, arc=3pt},
  nota/.style={enhanced, colback=rojo!5, colframe=rojo!40,
    top=4pt, bottom=4pt, left=6pt, right=6pt, arc=3pt, fontupper=\small},
}

\DeclareMathOperator{\Var}{Var}
\DeclareMathOperator{\Cov}{Cov}
\newcommand{\WN}{\mathrm{WN}}
\newcommand{\AIC}{\mathrm{AIC}}
\newcommand{\BIC}{\mathrm{BIC}}

\begin{document}

%% ── PORTADA ─────────────────────────────────────────────────────────────────
\begin{titlepage}
\centering\vspace*{1.5cm}
{\Large\bfseries\color{navy} UNIVERSIDAD PANAMERICANA\par}
\vspace{0.2cm}{\normalsize Facultad de Ciencias Empresariales\par}
\vspace{1.2cm}\textcolor{navy}{\rule{\linewidth}{2pt}}\vspace{0.4cm}
{\Huge\bfseries\color{navy} Modelos \textcolor{rojo}{SARIMA}\par
  con Datos Financieros\par}
\vspace{0.3cm}{\large Forma General, Interpretaci\'on y Tres Ejemplos en Python\par}
\vspace{0.4cm}\textcolor{navy}{\rule{\linewidth}{2pt}}\vspace{1.2cm}
\begin{tcolorbox}[formula, width=0.94\linewidth]
\centering\textbf{Ecuaci\'on General SARIMA$(p,d,q)(P,D,Q)_s$}\\[8pt]
\large
$\underbrace{\Phi_P(B^s)}_{\text{AR estac.}}
 \underbrace{\varphi_p(B)}_{\text{AR ord.}}
 \underbrace{(1-B^s)^D}_{\Delta_s^D}
 \underbrace{(1-B)^d}_{\Delta^d}
 Y_t
 \;=\;
 \underbrace{\Theta_Q(B^s)}_{\text{MA estac.}}
 \underbrace{\theta_q(B)}_{\text{MA ord.}}
 \underbrace{\varepsilon_t}_{\WN(0,\sigma^2)}$
\end{tcolorbox}
\vspace{1.5cm}
\begin{tabular}{rl}
  \textbf{Materia:}   & Series de Tiempo para Finanzas \\[4pt]
  \textbf{Profesor:}  & Dr.\ C\'esar Galindo \\[4pt]
  \textbf{Per\'iodo:} & Verano 2026 \\[4pt]
  \textbf{Fecha:}     & \today \\
\end{tabular}
\vfill
{\small\color{gray}
 Datos: \texttt{statsmodels} (T-Bills, CO$_2$ Mauna Loa, PIB Real EE.UU.) $\cdot$
 C\'odigo: Python~3 / \texttt{statsmodels.tsa.statespace.SARIMAX}}
\end{titlepage}

\tableofcontents\newpage

%% ── SECCION 1: FORMA GENERAL ─────────────────────────────────────────────────
\section{Forma General del Modelo SARIMA$(p,d,q)(P,D,Q)_s$}

\subsection{Ecuaci\'on completa}

El modelo SARIMA combina una parte no estacional ARIMA$(p,d,q)$ con una parte
estacional ARIMA$(P,D,Q)_s$:

\begin{tcolorbox}[formula, title={Ecuaci\'on General}]
\begin{equation}
  \Phi_P(B^s)\;\varphi_p(B)\;(1-B^s)^D\;(1-B)^d\;Y_t
  \;=\;
  \Theta_Q(B^s)\;\theta_q(B)\;\varepsilon_t,
  \quad\varepsilon_t\sim\WN(0,\sigma^2)
  \label{eq:sarima}
\end{equation}
con los operadores:
\begin{align*}
  \varphi_p(B)  &= 1-\varphi_1 B-\cdots-\varphi_p B^p
  &&\text{(AR ordinario)}\\
  \theta_q(B)   &= 1+\theta_1 B+\cdots+\theta_q B^q
  &&\text{(MA ordinario)}\\
  \Phi_P(B^s)   &= 1-\Phi_1 B^s-\cdots-\Phi_P B^{Ps}
  &&\text{(AR estacional)}\\
  \Theta_Q(B^s) &= 1+\Theta_1 B^s+\cdots+\Theta_Q B^{Qs}
  &&\text{(MA estacional)}
\end{align*}
\end{tcolorbox}

\subsection{Los siete par\'ametros}

\begin{table}[H]
\centering
\caption{Par\'ametros del modelo SARIMA$(p,d,q)(P,D,Q)_s$}
\renewcommand{\arraystretch}{1.35}
\begin{tabular}{>{\bfseries\color{rojo}\centering\arraybackslash}m{1cm}
                >{\centering\arraybackslash}m{2.3cm}m{5.5cm}m{4cm}}
\toprule
\rowcolor{navy!12}
\color{navy}Par\'am. & \color{navy}Tipo &
\color{navy}Descripci\'on & \color{navy}Identificaci\'on\\
\midrule
$p$ & AR ordinario   & Rezagos de $Y_t$ en la ecuaci\'on & PACF corte en lag $p$ \\
$d$ & Integ.\ ord.   & Diferencias $(1-B)^d$ & Prueba ADF / KPSS \\
$q$ & MA ordinario   & Errores pasados $\varepsilon_{t-k}$ & ACF corte en lag $q$ \\
\midrule
\rowcolor{grisclaro}
$P$ & AR estacional  & Rezagos $Y_{t-ks}$, $k=1,\ldots,P$ & PACF pico en $s,2s,\ldots$ \\
\rowcolor{grisclaro}
$D$ & Integ.\ estac. & Diferencias $(1-B^s)^D$ & Gr\'afica / prueba CH \\
\rowcolor{grisclaro}
$Q$ & MA estacional  & Errores en m\'ultiplos de $s$ & ACF pico en $s,2s,\ldots$ \\
\midrule
$s$ & Per\'iodo & Longitud ciclo (12, 4, 52\ldots) & Inspecci\'on visual \\
\bottomrule
\end{tabular}
\end{table}

\subsection{Reglas de identificaci\'on ACF/PACF}

\begin{table}[H]
\centering
\caption{Patrones para identificar el orden del modelo}
\renewcommand{\arraystretch}{1.3}
\begin{tabular}{lll}
\toprule
\rowcolor{navy!12}
\color{navy}\textbf{Proceso} & \color{navy}\textbf{ACF} & \color{navy}\textbf{PACF}\\
\midrule
AR$(p)$     & Decae geom\'etricamente & Se corta en rezago $p$ \\
MA$(q)$     & Se corta en rezago $q$ & Decae geom\'etricamente \\
ARMA$(p,q)$ & Decae (mezcla) & Decae (mezcla) \\
\rowcolor{grisclaro}
SAR$(P)_s$  & Picos en $s,2s,\ldots$ decaen & Pico significativo en $s$ \\
\rowcolor{grisclaro}
SMA$(Q)_s$  & Pico significativo en $s$ & Picos en $s,2s,\ldots$ decaen \\
\bottomrule
\end{tabular}
\end{table}

\subsection{Criterios de selecci\'on y diagn\'ostico}

\begin{tcolorbox}[nota]
$\AIC = -2\hat\ell + 2k$ \quad $\BIC = -2\hat\ell + k\ln n$ \quad
Minimizar ambos. Con $n>100$ usar \textbf{BIC}. Parsimonia: $p+d+q+P+D+Q\le 6$.
\end{tcolorbox}

\begin{tcolorbox}[nota]
\textbf{Ljung-Box:}
$Q_{LB} = n(n+2)\sum_{k=1}^{h}\hat\rho_k^2/(n-k)\sim\chi^2(h-p-q)$ bajo $H_0$.
Si $p\text{-valor}>0.05 \Rightarrow$ residuos son ruido blanco \checkmark.
\end{tcolorbox}

\begin{figure}[H]\centering
\includegraphics[width=\linewidth]{ej0_forma_general_sarima.png}
\caption{Forma general SARIMA y tabla de interpretaci\'on de par\'ametros.}
\end{figure}

\newpage

%% ── SECCION 2: EJEMPLO 1 ─────────────────────────────────────────────────────
\section{Ejemplo 1 --- Tasa de Inter\'es Trimestral (proxy CETES/T-Bills)}

\begin{tcolorbox}[teorema, title={Modelo: SARIMA$(1,1,0)(1,0,0)_4$}]
\textbf{Serie:} Tasa T-Bill trimestral EE.UU.\ 1959--2009 (203 obs.), proxy CETES.\\[4pt]
\textbf{Ecuaci\'on ajustada:}
\[
  (1-\hat\varphi_1 B)(1-\hat\Phi_1 B^4)(1-B)\,Y_t = \varepsilon_t
\]
\end{tcolorbox}

\subsection{Descripci\'on y prueba de ra\'iz unitaria}

Media $5.31\%$, desviaci\'on est\'andar $2.80\%$.
ADF: estad\'istico $-2.04$, $p=0.270 \Rightarrow$ I(1), $d=1$.
PACF de $\Delta Y_t$ corte en lag~1 ($p=1$), pico estacional en lag~4 ($P=1$).

\subsection{Resultados de estimaci\'on}

\begin{table}[H]\centering
\caption{Estimaci\'on SARIMA$(1,1,0)(1,0,0)_4$ --- Tasa T-Bill}
\renewcommand{\arraystretch}{1.3}
\begin{tabular}{lrrrr}
\toprule
\rowcolor{navy!12}\color{navy}Par\'ametro &
\color{navy}Estimado & \color{navy}Err.\ est. &
\color{navy}$z$ & \color{navy}$p$-valor\\
\midrule
Intercepto     & $-0.013$ & $0.070$ & $-0.18$ & $0.859$\\
$\hat\varphi_1$& $\phantom{-}0.052$ & $0.032$ & $\phantom{-}1.64$ & $0.101$\\
$\hat\Phi_1$   & $-0.045$ & $0.047$ & $-0.96$ & $0.340$\\
$\hat\sigma^2$ & $\phantom{-}0.761$ & $0.029$ & $25.83$ & $<0.001$\\
\midrule
\multicolumn{2}{l}{$\AIC=513.34$} & \multicolumn{3}{l}{$\BIC=526.47$}\\
\bottomrule
\end{tabular}
\end{table}

\begin{tcolorbox}[aplicacion, title={Aplicaci\'on: Valuaci\'on de Bonos / CETES}]
El pron\'ostico de tasas a 8 trimestres permite calcular el valor presente de
bonos de tasa fija, decidir coberturas fija/variable (swaps) y estimar duraci\'on modificada.
\end{tcolorbox}

\begin{figure}[H]\centering
\includegraphics[width=\linewidth]{ej1_tasa_sarima.png}
\caption{SARIMA$(1,1,0)(1,0,0)_4$ --- Serie original, ACF/PACF de $\Delta Y_t$,
         residuos y pron\'ostico 8 trimestres con IC~95\%.}
\end{figure}

\newpage

%% ── SECCION 3: EJEMPLO 2 ─────────────────────────────────────────────────────
\section{Ejemplo 2 --- CO\textsubscript{2} Mensual: Modelo de Aerolineas (Airline)}

\begin{tcolorbox}[teorema, title={Modelo: SARIMA$(0,1,1)(0,1,1)_{12}$ --- Airline de Box-Jenkins}]
\textbf{Serie:} CO$_2$ mensual Mauna Loa 1958--2001 (526 obs.),
$W_t=\ln Y_t$ para estabilizar varianza.\\[4pt]
\textbf{Ecuaci\'on:}
$(1-B)(1-B^{12})\,W_t = (1+\hat\theta_1 B)(1+\hat\Theta_1 B^{12})\,\varepsilon_t$
\end{tcolorbox}

\subsection{Expansi\'on del operador (Ejercicio 3.2 de la tarea)}

\begin{align}
  (1-B)(1-B^{12})\,W_t &= W_t - W_{t-1} - W_{t-12} + W_{t-13}\label{eq:lhs}\\
  (1+\theta_1 B)(1+\Theta_1 B^{12})\,\varepsilon_t
  &= \varepsilon_t + \theta_1\varepsilon_{t-1}
     + \Theta_1\varepsilon_{t-12} + \theta_1\Theta_1\varepsilon_{t-13}\label{eq:rhs}
\end{align}

\subsection{Resultados}

\begin{table}[H]\centering
\caption{Estimaci\'on SARIMA$(0,1,1)(0,1,1)_{12}$ --- CO$_2$}
\renewcommand{\arraystretch}{1.3}
\begin{tabular}{lrrrr}
\toprule
\rowcolor{navy!12}\color{navy}Par\'ametro &
\color{navy}Estimado & \color{navy}Err.\ est. &
\color{navy}$z$ & \color{navy}$p$-valor\\
\midrule
$\hat\theta_1$ (MA L1)   & $-0.3418$ & $0.038$ & $-8.93$ & $<0.001$\\
$\hat\Theta_1$ (SMA L12) & $-0.6798$ & $0.041$ & $-16.45$ & $<0.001$\\
$\hat\sigma^2$ & $8.29\!\times\!10^{-7}$ & --- & $15.41$ & $<0.001$\\
\midrule
\multicolumn{2}{l}{$\AIC=-5562.55$} & \multicolumn{3}{l}{$\BIC=-5549.91$ \quad\textbf{Ganador}}\\
\bottomrule
\end{tabular}
\end{table}

\begin{tcolorbox}[aplicacion, title={Aplicaci\'on: Commodities Energ\'eticos}]
Est\'andar en gas natural (ciclo invernal), electricidad, demanda de energ\'ia.
Permite construir \emph{hedge ratios} estacionales y fijar precios de futuros.
\end{tcolorbox}

\begin{figure}[H]\centering
\includegraphics[width=\linewidth]{ej2_co2_airline.png}
\caption{SARIMA$(0,1,1)(0,1,1)_{12}$ --- Transformaci\'on log, $\Delta_1\Delta_{12}$,
         ACF estacional, Q-Q de residuos y pron\'ostico 24 meses.}
\end{figure}

\newpage

%% ── SECCION 4: EJEMPLO 3 ─────────────────────────────────────────────────────
\section{Ejemplo 3 --- PIB Real Trimestral: Selecci\'on Autom\'atica Box-Jenkins}

\begin{tcolorbox}[teorema, title={Modelo ganador por BIC: SARIMA$(1,1,1)(1,0,0)_4$}]
\textbf{Serie:} PIB real EE.UU.\ trimestral 1959--2009 (203 obs., en log).
Crecimiento promedio: $3.14\%$ anual.\\[4pt]
\textbf{Ecuaci\'on ajustada:}
\[
  (1-0.923\,B)(1-0.018\,B^4)(1-B)\,\ln(\text{PIB})_t
  = (1-0.567\,B)\,\varepsilon_t,\quad\hat\sigma=0.00843
\]
\end{tcolorbox}

\subsection{Metodolog\'ia Box-Jenkins}

\subsubsection{Identificaci\'on}
ADF: $p=0.999 \Rightarrow$ I(1), $d=1$. ACF de $\Delta\ln\text{PIB}$ decae; PACF corte en lag~1.
\subsubsection{Estimaci\'on en grilla}

\begin{table}[H]\centering
\caption{Top~5 modelos por BIC --- PIB trimestral}
\renewcommand{\arraystretch}{1.3}
\begin{tabular}{lrrr}
\toprule
\rowcolor{navy!12}\color{navy}Modelo &
\color{navy}AIC & \color{navy}BIC & \color{navy}Params.\\
\midrule
\rowcolor{verde!10}SARIMA$(1,1,1)(0,0,0)_4$ & $-1331.5$ & $-1321.6$ & 3\\
SARIMA$(2,1,0)(0,0,0)_4$ & $-1325.3$ & $-1315.4$ & 3\\
SARIMA$(1,1,2)(0,0,0)_4$ & $-1328.0$ & $-1314.8$ & 4\\
SARIMA$(2,1,1)(0,0,0)_4$ & $-1327.2$ & $-1314.0$ & 4\\
SARIMA$(2,1,2)(0,0,0)_4$ & $-1319.9$ & $-1303.5$ & 5\\
\bottomrule
\end{tabular}
\end{table}

\subsubsection{ECM del pron\'ostico (fan chart)}

\[
  \mathrm{ECM}_h = \sigma^2\sum_{j=0}^{h-1}\psi_j^2,
  \quad \psi_0=1,\quad \psi_j=\varphi_1\psi_{j-1}+\varphi_2\psi_{j-2}+\cdots
\]

El fan chart se expande porque $\mathrm{ECM}_h$ crece mon\'otonamente con $h$.

\begin{table}[H]\centering
\caption{Pron\'ostico PIB real (miles de MM USD)}
\renewcommand{\arraystretch}{1.3}
\begin{tabular}{crrr}
\toprule
\rowcolor{navy!12}\color{navy}Trimestre &
\color{navy}Pron\'ostico & \color{navy}LI 95\% & \color{navy}LS 95\%\\
\midrule
2009-Q4 & $12{,}976.6$ & $12{,}763.9$ & $13{,}192.8$\\
2010-Q1 & $12{,}962.9$ & $12{,}607.0$ & $13{,}329.0$\\
2010-Q2 & $12{,}953.6$ & $12{,}453.4$ & $13{,}473.9$\\
2010-Q3 & $12{,}947.0$ & $12{,}299.6$ & $13{,}628.5$\\
2010-Q4 & $12{,}939.2$ & $12{,}140.3$ & $13{,}790.6$\\
2011-Q1 & $12{,}931.9$ & $11{,}980.3$ & $13{,}959.2$\\
\bottomrule
\end{tabular}
\end{table}

\begin{tcolorbox}[aplicacion, title={Aplicaci\'on: Riesgo Crediticio y Ciclos de Negocios}]
Calibrar modelos de riesgo crediticio (PD, LGD) bajo escenarios macro;
asignaci\'on t\'actica de activos; valuaci\'on de empresas c\'iclicas.
\end{tcolorbox}

\begin{figure}[H]\centering
\includegraphics[width=\linewidth]{ej3_pib_boxjenkins.png}
\caption{SARIMA$(1,1,1)(1,0,0)_4$ --- Crecimiento trimestral, comparaci\'on por BIC,
         ACF/PACF y fan chart IC 68\%--95\% a 12 trimestres.}
\end{figure}

\newpage

%% ── SECCION 5: RESUMEN ───────────────────────────────────────────────────────
\section{Resumen: Gu\'ia R\'apida para el Examen}

\begin{tcolorbox}[formula, title={Checklist Box-Jenkins}]
\begin{enumerate}[noitemsep]
  \item \textbf{Graficar} la serie. \textquestiondown Tendencia? \textquestiondown Estacionalidad?
  \item \textbf{ADF/KPSS:} $p>0.05 \Rightarrow d=1$; picos ACF en $ks \Rightarrow D=1$.
  \item \textbf{ACF/PACF} diferenciada: PACF corte $\to p$; ACF corte $\to q$; picos en $s,2s\ldots \to P,Q$.
  \item \textbf{Estimar} candidatos $\to$ calcular AIC/BIC $\to$ elegir menor BIC.
  \item \textbf{Ljung-Box:} $p>0.05 \Rightarrow$ modelo OK \checkmark.
  \item \textbf{Pron\'ostico} con intervalos que crecen con $h$.
\end{enumerate}
\end{tcolorbox}

\begin{table}[H]\centering
\caption{Resumen de los tres ejemplos}
\renewcommand{\arraystretch}{1.35}
\begin{tabular}{llllll}
\toprule
\rowcolor{navy!12}
\color{navy} & \color{navy}Serie & \color{navy}Modelo &
\color{navy}AIC & \color{navy}BIC & \color{navy}Aplicaci\'on\\
\midrule
\textbf{Ej.1} & Tasa T-Bill trim.   & SARIMA$(1,1,0)(1,0,0)_4$    & $513.3$   & $526.5$   & Bonos/CETES\\
\textbf{Ej.2} & CO$_2$ mensual      & SARIMA$(0,1,1)(0,1,1)_{12}$ & $-5562.6$ & $-5549.9$ & Commodities\\
\textbf{Ej.3} & PIB real trim.      & SARIMA$(1,1,1)(1,0,0)_4$    & $-1314.1$ & $-1301.0$ & Riesgo cred.\\
\bottomrule
\end{tabular}
\end{table}

\newpage

%% ── SECCION 6: CODIGO ────────────────────────────────────────────────────────
\section{C\'odigo Python Completo}

\begin{tcolorbox}[nota]
\textbf{Dependencias:} \texttt{numpy}, \texttt{pandas}, \texttt{matplotlib},
\texttt{statsmodels}, \texttt{scipy}. Sin necesidad de conexi\'on a internet.
Ejecutar: \texttt{python3 sarima\_clean.py}
\end{tcolorbox}

\lstinputlisting[
  style=python,
  caption={C\'odigo completo: tres ejemplos SARIMA con datos financieros},
  label={lst:sarima},
]{sarima_clean.py}

%% ── REFERENCIAS ──────────────────────────────────────────────────────────────
\newpage
\section*{Referencias}
\addcontentsline{toc}{section}{Referencias}
\begin{enumerate}[label={[\arabic*]}]
  \item Box, G.E.P.\ y Jenkins, G.M.\ (1976). \textit{Time Series Analysis}. Holden-Day.
  \item Hamilton, J.D.\ (1994). \textit{Time Series Analysis}. Princeton University Press.
  \item Shumway, R.H.\ y Stoffer, D.S.\ (2017). \textit{Time Series Analysis and Its Applications}. Springer.
  \item Hyndman, R.J.\ y Athanasopoulos, G.\ (2021). \textit{Forecasting: Principles and Practice}. \url{https://otexts.com/fpp3/}
  \item Statsmodels Team (2024). \url{https://www.statsmodels.org}
\end{enumerate}

\end{document}
